\subsection{5. Phương pháp nghiên cứu}

\indent\indent Đề tài triển khai kết hợp các phương pháp nghiên cứu sau:\vspace{0.3cm}

\begin{enumerate}[label=\alph*)]
    \item \textbf{Phương pháp nghiên cứu lý thuyết:} Khảo sát toán học về phép biến đổi DCT-II trực chuẩn trên khối $8 \times 8$, nguyên lý Compound Scaling của EfficientNet, cơ chế chú ý scSE và cấu trúc giải mã lặp đa góc nhìn MID.
    \item \textbf{Phương pháp tiền xử lý và kỹ thuật dữ liệu:} Lập trình thuật toán bóc tách nhãn từ văn bản chú thích, xử lý cụm từ phủ định, tái cân bằng tập dữ liệu và áp dụng pipeline tăng cường ảnh (Data Augmentation).
    \item \textbf{Phương pháp thực nghiệm mô hình hóa:}
    \begin{itemize}
        \item Cài đặt mô-đun BlockDCT tính toán nhanh trên PyTorch tensor.
        \item Xây dựng nhánh kết hợp không gian -- tần số tại độ phân giải $1/8$.
        \item Thiết kế và cài đặt bộ giải mã lồng nhau LiteMID sử dụng Depthwise Separable Convolutions.
        \item Xây dựng hàm mất mát kết hợp `MultiTaskTamperLoss` gồm BCE + Dice Loss cho nhánh phân đoạn và BCE cho nhánh phân loại.
    \end{itemize}
    \item \textbf{Phương pháp đánh giá thực nghiệm:} Đánh giá trên tập kiểm thử độc lập qua các chỉ số Pixel Accuracy, Precision, Recall, IoU, Macro Dice và Top-3 Class Hit Rate.
\end{enumerate}
